\chapter{Nauwere verweving tussen code \& documentatie}
\label{ch:verweving_code_docs}

De uiteindelijke meerwaarde van de voorgestelde aanpak ligt in de fundamentele verschuiving van documentatie als een statisch artefact naar documentatie als een integraal onderdeel van de software-levenscyclus. Door de scheiding tussen deze twee werelden op te heffen, ontstaat een ecosysteem dat zowel wetenschappelijk accuraat als technisch robuust is.

\section{Minder communicatie-overhead en ruis}
Een van de meest tastbare voordelen van de nauwere verweving is de drastische daling in de communicatie-overhead. In de traditionele workflow ging veel tijd verloren aan het verduidelijken van specificaties. Onderzoekers moesten wachten op IT om te zien hoe een formule in de praktijk presteerde, waarna IT vaak moest terugkomen bij de onderzoekers omdat de wiskundige beschrijving in het Word-document ambigu was.

In de nieuwe workflow fungeert de Jupyter-notebook als een gemeenschappelijke taal (ubiquitous language). De programmeur ziet de C\#-interop en de achterliggende types, terwijl de onderzoeker de resultaten en de narratieve tekst ziet. Dit gedeelde referentiekader minimaliseert de kans op spraakverwarring. Vragen over de interne werking van de code kunnen nu vaak in enkele minuten door de onderzoeker zelf worden beantwoord door simpelweg een code-cel in de Jupyter-notebook uit te voeren.

\section{Documentatie als \textquotedbl Single Source of Truth\textquotedbl}
Door de principes van Literate Programming te omarmen, wordt documentatie niet langer een sluitpost aan het einde van een project, maar de drijvende kracht erachter. Elke wijziging in een wetenschappelijk model wordt gelijktijdig gedocumenteerd (in tekst) en geïmplementeerd (in code) binnen hetzelfde proces.

Dit verhoogt de kwaliteit van het intellectueel eigendom van ILVO aanzienlijk. De kennis zit niet meer gefragmenteerd in de hoofden van individuele experts of in een wirwar van ongecontroleerde scripts op lokale computers. In plaats daarvan is er één traceerbare bron waar de wetenschappelijke redenering en de technische executie samenkomen. Dit is essentieel voor de reproduceerbaarheid van het onderzoek, aangezien men jaren later nog exact kan zien welke aannames tot welke codewijzigingen hebben geleid.

\section{Conclusie: een strategische transformatie}
De integratie van geavanceerde software-architectuur, container-isolatie en interactieve documentatievormen transformeert de relatie tussen IT en onderzoek fundamenteel. De IT-afdeling verschuift van een uitvoerende partij naar een platform-enabler. Zij faciliteren de infrastructuur en de robuuste backend, terwijl de onderzoekers de domeinlogica voeden, verfijnen en valideren met een voorheen ongekende mate van autonomie.

Deze nauwe verweving vormt een veelbelovende oplossingsrichting voor de complexiteitsuitdagingen binnen EMAV-web. Het onderzoek en de bijbehorende proof-of-concept tonen aan dat een platform waarop technologische vernieuwing en wetenschappelijke precisie elkaar versterken, technisch haalbaar is. Door de kloof tussen code en documentatie te verkleinen, wordt een fundament gelegd dat de potentie heeft om de wendbaarheid en de kwaliteit van het wetenschappelijk onderzoek bij ILVO op lange termijn structureel te ondersteunen. De succesvolle validatie van dit concept opent de deur naar een meer geïntegreerde werkwijze waarbij IT en wetenschap elkaar optimaal versterken.
